% !PRINTED 442
\weberpage{442}
% running head: 442  Eighteenth Chapter.  § 122.

%<<CONT
\gesperrt{and since $\mathfrak{S}$ is an Abelian group, the class equation is an Abelian one.}

Let us assume that the function $R$ has real (rational) coefficients; its rational expression will nevertheless have the form $a+b\sqrt{D}$, in which $a, b$ are rational numbers. The imaginary part will drop out then and only then, when $R$ also remains unchanged under the interchange of all the classes $k, k', k'', \ldots$ with their opposites. From this it follows that, without adjunction of $\sqrt{D}$, the \gesperrt{Galois} group of the class equation is contained in the system of $2h$ permutations which one obtains by converting every class in $\mathfrak{S}$ into its opposite. Only in the particular case in which all the classes are two-sided (which occurs only for a finite number of determinants) is this latter system identical with $\mathfrak{S}$, and the class equation is then an \gesperrt{Abelian} one without adjunction of $\sqrt{D}$.

The algebraic field which is formed from the rational functions of a class invariant is therefore, apart from the exceptional case last mentioned, not a \gesperrt{normal field}, but is different from its conjugate fields. One obtains a normal field, on the other hand, if one adjoins the square root $\sqrt{D}$ to the field of the class invariants. For every rational function of all the roots of the equation can be expressed rationally by one of them and $\sqrt{D}$.

\begin{center}
§ 122. \gesperrt{Irreducibility.}
\end{center}

We now consider the algebraic number field
\weq{1}{\mathfrak{K} = \mathfrak{R}(k, \sqrt{D}),}
which is composed of a class invariant $(k)$ of the discriminant $D$ and of $\sqrt{D}$; by the preceding section this field contains all the class invariants $(k), (k'), (k''), \ldots$ belonging to the same discriminant. This field we call the \gesperrt{class field} of the discriminant $D$. The quadratic field $\mathfrak{R}(\sqrt{D})$ we denote, as hitherto, by $\Omega$; if $D=Q^2\Delta$ and $\Delta$ is the stem of $D$, it is identical with $\mathfrak{R}(\sqrt{\Delta})$.
Let $h$ be the \gesperrt{class number} of the discriminant $D$.

1. Let $(k)$ be any one of the $h$ class variants, and let % [sic] German print has "Klassenvarianten"; context requires "Klasseninvarianten"
\weq{2}{\Phi(k) = (k)^\nu + \alpha_1(k)^{\nu-1} + \alpha_2(k)^{\nu-2}\cdots + \alpha_\nu = 0}
be the equation of lowest degree in $\Omega$ whose root is $(k)$. Since $(k)$ is an integer, the $\alpha_i$ must also be integers
%CONT>>
